% GENERATED FILE. Edit canonical lessons, then run npm run generate:books.
% canonical-source-hash: fnv1a64:23dbc9375f26683e
% canonical-lessons: GU-C118-masalo, GU-C118-haldar, GU-C118-jiru, GU-C118-ghas, GU-C118-tapu, GU-R118-first-pass-the-house-and-clothes, GU-R118-second-pass-clothes-and-food

\chapter{Spice, Turmeric, Cumin, Grass, Island}
\label{ch:gu-spice-turmeric-cumin-grass-island}
\hlchaptermodality{118}

\begin{chapteropening}
\textbf{By the end of this chapter:} \emph{I can say spice, turmeric, cumin, grass and an island.}

Complete the last lesson of chapter 118: I can say spice, turmeric, cumin, grass and an island.
\end{chapteropening}

\section[\texorpdfstring{\gu{મસાલો} (masālo)}{મસાલો (masālo)}]{\gu{મસાલો} (masālo) — spice}

\label{lesson:GU-C118-masalo}



\begin{quote}
Before the new one: say the Gujarati for halwa, then the Gujarati for rice pudding.

Between two sentences \emph{paṇ} means \textquotedblleft{}but.\textquotedblright{} What does it mean right after a noun? (\textquotedblleft{}Also.\textquotedblright{})
\end{quote}

\subsection*{You'll want to know: \gu{મસાલો}}
\textbf{\gu{મસાલો}} (\emph{masālo}) — \textquotedblleft{}spice\textquotedblright{}.

\begin{grammarlens}[title={What you've built}]
One more word for this chapter.
\end{grammarlens}

\subsection*{Your turn}
\begin{itemize}
\raggedright
  \item \emph{Say it:} \emph{masālo}
  \item \emph{Say it:} \emph{masālo}, once more
  \item \emph{Say it:} say \emph{masālo}
  \item \emph{Recall:} say the Gujarati for halwa, then the Gujarati for rice pudding, then say \emph{masālo} again
\end{itemize}

\begin{tcolorbox}[breakable,colback=teal!4,colframe=teal!35!black,title={Before you move on}]
What does \gu{મસાલો} mean? (\textquotedblleft{}Spice\textquotedblright{}.) Say it once more.
\end{tcolorbox}
\section[\texorpdfstring{\gu{હળદર} (haḷdar)}{હળદર (haḷdar)}]{\gu{હળદર} (haḷdar) — turmeric}

\label{lesson:GU-C118-haldar}



\begin{quote}
Before the new one: say the Gujarati for rice pudding, then the Gujarati for spice.
\end{quote}

\subsection*{You'll want to know: \gu{હળદર}}
\textbf{\gu{હળદર}} (\emph{haḷdar}) — \textquotedblleft{}turmeric\textquotedblright{}.

\begin{grammarlens}[title={What you've built}]
One more word for this chapter.
\end{grammarlens}

\subsection*{Your turn}
\begin{itemize}
\raggedright
  \item \emph{Say it:} \emph{haḷdar}
  \item \emph{Say it:} \emph{haḷdar}, once more
  \item \emph{Say it:} say \emph{haḷdar}
  \item \emph{Recall:} say the Gujarati for rice pudding, then the Gujarati for spice, then say \emph{haḷdar} again
\end{itemize}

\begin{tcolorbox}[breakable,colback=teal!4,colframe=teal!35!black,title={Before you move on}]
What does \gu{હળદર} mean? (\textquotedblleft{}Turmeric\textquotedblright{}.) Say it once more.
\end{tcolorbox}
\section[\texorpdfstring{\gu{જીરું} (jīrũ)}{જીરું (jīrũ)}]{\gu{જીરું} (jīrũ) — cumin}

\label{lesson:GU-C118-jiru}



\begin{quote}
Before the new one: say the Gujarati for spice, then the Gujarati for turmeric.
\end{quote}

\subsection*{You'll want to know: \gu{જીરું}}
\textbf{\gu{જીરું}} (\emph{jīrũ}) — \textquotedblleft{}cumin\textquotedblright{}.

\begin{grammarlens}[title={What you've built}]
One more word for this chapter.
\end{grammarlens}

\subsection*{Your turn}
\begin{itemize}
\raggedright
  \item \emph{Say it:} \emph{jīrũ}
  \item \emph{Say it:} \emph{jīrũ}, once more
  \item \emph{Say it:} say \emph{jīrũ}
  \item \emph{Recall:} say the Gujarati for spice, then the Gujarati for turmeric, then say \emph{jīrũ} again
\end{itemize}

\begin{tcolorbox}[breakable,colback=teal!4,colframe=teal!35!black,title={Before you move on}]
What does \gu{જીરું} mean? (\textquotedblleft{}Cumin\textquotedblright{}.) Say it once more.
\end{tcolorbox}
\section[\texorpdfstring{\gu{ઘાસ} (ghās)}{ઘાસ (ghās)}]{\gu{ઘાસ} (ghās) — grass}

\label{lesson:GU-C118-ghas}



\begin{quote}
Before the new one: say the Gujarati for turmeric, then the Gujarati for cumin.
\end{quote}

\subsection*{You'll want to know: \gu{ઘાસ}}
\textbf{\gu{ઘાસ}} (\emph{ghās}) — \textquotedblleft{}grass\textquotedblright{}.

\begin{grammarlens}[title={What you've built}]
One more word for this chapter.
\end{grammarlens}

\subsection*{Your turn}
\begin{itemize}
\raggedright
  \item \emph{Say it:} \emph{ghās}
  \item \emph{Say it:} \emph{ghās}, once more
  \item \emph{Say it:} say \emph{ghās}
  \item \emph{Recall:} say the Gujarati for turmeric, then the Gujarati for cumin, then say \emph{ghās} again
\end{itemize}

\begin{tcolorbox}[breakable,colback=teal!4,colframe=teal!35!black,title={Before you move on}]
What does \gu{ઘાસ} mean? (\textquotedblleft{}Grass\textquotedblright{}.) Say it once more.
\end{tcolorbox}
\section[\texorpdfstring{\gu{ટાપુ} (ṭāpu)}{ટાપુ (ṭāpu)}]{\gu{ટાપુ} (ṭāpu) — an island}

\label{lesson:GU-C118-tapu}



\begin{quote}
Before the new one: say the Gujarati for cumin, then the Gujarati for grass.
\end{quote}

\subsection*{You'll want to know: \gu{ટાપુ}}
\textbf{\gu{ટાપુ}} (\emph{ṭāpu}) — \textquotedblleft{}an island\textquotedblright{}.

\begin{grammarlens}[title={What you've built}]
One more word for this chapter.
\end{grammarlens}

\subsection*{Your turn}
\begin{itemize}
\raggedright
  \item \emph{Say it:} \emph{ṭāpu}
  \item \emph{Say it:} \emph{ṭāpu}, once more
  \item \emph{Say it:} say \emph{ṭāpu}
  \item \emph{Recall:} say the Gujarati for cumin, then the Gujarati for grass, then say \emph{ṭāpu} again
\end{itemize}

\begin{tcolorbox}[breakable,colback=teal!4,colframe=teal!35!black,title={Before you move on}]
What does \gu{ટાપુ} mean? (\textquotedblleft{}An island\textquotedblright{}.) Say it once more.
\end{tcolorbox}
\section[(dialogue)]{The house and clothes — a first pass}

\label{lesson:GU-R118-first-pass-the-house-and-clothes}



\begin{quote}
Say the words for grass and an island.
\end{quote}

\subsection*{Your turn}
Say these aloud:

\begin{itemize}
\raggedright
  \item a box, a jar, a candle, a pillow, a doll
  \item a kurta, pyjamas, a ring, an earring, socks
  \item shoes, a belt, a handkerchief, a sweater, a coat
  \item a pocket, jaggery, honey, chickpeas, mung beans
\end{itemize}

\begin{tcolorbox}[breakable,colback=teal!4,colframe=teal!35!black,title={Before you move on}]
A box? (\textbf{\gu{પેટી}}.) A jar? (\textbf{\gu{બરણી}}.) A candle? (\textbf{\gu{મીણબત્તી}}.) A pillow? (\textbf{\gu{તકિયો}}.) A doll? (\textbf{\gu{ઢીંગલી}}.) A kurta? (\textbf{\gu{કુરતો}}.) Pyjamas? (\textbf{\gu{પાયજામો}}.) A ring? (\textbf{\gu{વીંટી}}.) An earring? (\textbf{\gu{બુટ્ટી}}.) Socks? (\textbf{\gu{મોજાં}}.) Shoes? (\textbf{\gu{બૂટ}}.) A belt? (\textbf{\gu{પટ્ટો}}.) A handkerchief? (\textbf{\gu{રૂમાલ}}.) A sweater? (\textbf{\gu{સ્વેટર}}.) A coat? (\textbf{\gu{કોટ}}.) A pocket? (\textbf{\gu{ખિસ્સું}}.) Jaggery? (\textbf{\gu{ગોળ}}.) Honey? (\textbf{\gu{મધ}}.) Chickpeas? (\textbf{\gu{ચણા}}.) Mung beans? (\textbf{\gu{મગ}}.)
\end{tcolorbox}
\section[(dialogue)]{Clothes and food — a second pass}

\label{lesson:GU-R118-second-pass-clothes-and-food}



\begin{quote}
Say the words for millet and wheat.
\end{quote}

\subsection*{Your turn}
Say these aloud:

\begin{itemize}
\raggedright
  \item millet, wheat, maize, an aubergine, cabbage
  \item a carrot, garlic, ginger, a lemon, a coconut
  \item a watermelon, a papaya, a guava, peanuts, dhokla
  \item kadhi, fritters, chutney, halwa, rice pudding
  \item spice, turmeric, cumin, grass, an island
\end{itemize}

\begin{tcolorbox}[breakable,colback=teal!4,colframe=teal!35!black,title={Before you move on}]
Millet? (\textbf{\gu{બાજરી}}.) Wheat? (\textbf{\gu{ઘઉં}}.) Maize? (\textbf{\gu{મકાઈ}}.) An aubergine? (\textbf{\gu{રીંગણ}}.) Cabbage? (\textbf{\gu{કોબી}}.) A carrot? (\textbf{\gu{ગાજર}}.) Garlic? (\textbf{\gu{લસણ}}.) Ginger? (\textbf{\gu{આદુ}}.) A lemon? (\textbf{\gu{લીંબુ}}.) A coconut? (\textbf{\gu{નાળિયેર}}.) A watermelon? (\textbf{\gu{તરબૂચ}}.) A papaya? (\textbf{\gu{પપૈયું}}.) A guava? (\textbf{\gu{જામફળ}}.) Peanuts? (\textbf{\gu{મગફળી}}.) Dhokla? (\textbf{\gu{ઢોકળાં}}.) Kadhi? (\textbf{\gu{કઢી}}.) Fritters? (\textbf{\gu{ભજિયાં}}.) Chutney? (\textbf{\gu{ચટણી}}.) Halwa? (\textbf{\gu{હલવો}}.) Rice pudding? (\textbf{\gu{ખીર}}.) Spice? (\textbf{\gu{મસાલો}}.) Turmeric? (\textbf{\gu{હળદર}}.) Cumin? (\textbf{\gu{જીરું}}.) Grass? (\textbf{\gu{ઘાસ}}.) An island? (\textbf{\gu{ટાપુ}}.)
\end{tcolorbox}
